\section{Environment selection}
\label{sec:env_selection}

The choice of environment is consequential because progress in RLlearning is closely coupled to the affordances and limitations of the simulator used. Classic platforms such as the Arcade Learning Environment enabled breakthroughs like DQN on Atari games, while more recent work on embodied and open-ended agents relies on visually rich 3D worlds such as Minecraft-based MineDojo. These environments expose a fundamental trade-off: researchers must often choose between computationally efficient but overly simplistic settings (e.g.\ small 2D gridworlds) and substantially slower but richer 3D simulators. Many of the richest environments are also built on closed-source commercial games such as Minecraft, which limits how much their mechanics can be modified or extended.

A second limitation is the lack of flexibility. Commonly used benchmarks either restrict customisation or require specifying environments in complex domain-specific languages, and many only support 2D worlds. This makes it difficult to design targeted scenarios for analysing specific behaviours such as long-horizon cooperation, coalition formation, or persistent social structure.

\paragraph{Why Craftium.} Craftium \cite{malagoncraftium} is a fully open-source 3D voxel world built on the Minetest game engine, supporting both single- and multi-agent scenarios through standard Gymnasium and PettingZoo interfaces. It is explicitly designed for rich 3D navigation, building, and resource-collection tasks, remains computationally efficient compared to Minecraft-based frameworks, and is highly customisable via a Lua API.\footnote{See \cite{malagoncraftium} for benchmarks and examples of single- and multi-agent tasks, open-world survival scenarios, and procedurally generated continual-learning curricula.} Its main drawback is that it is a framework rather than a fixed benchmark: additional engineering effort is required to design, implement, and validate specific social scenarios.

\paragraph{From OpenWorld to Five Chambers.} Our initial design used Craftium's OpenWorld environment with three heterogeneous roles (Engineer, Hunter, Guardian) corresponding to the three skill tracks (tools, hunt, defend). Pilot runs revealed that at small population sizes and short horizons, agents rarely co-located in shared regions, so co-activity signals required by the Hebbian update rule did not fire often enough for bonds to form. We considered scaling up the population and episode length to address this, but the compute requirements quickly exceeded the computation available for this research. We therefore designed \emph{Five Chambers}: a curriculum that retains Craftium's voxel-world affordances (3D navigation, resource collection, combat) but enforces co-location and shared task progression at every stage, so that the Hebbian mechanism can be evaluated under conditions where co-activity is reliably present. This environment structure now introduces more opportunities for collaboration, reducing the need for a longer horizon. Studying the mechanism at scale in genuinely open environments remains an open direction.

% \section{Experiment log}
% \label{sec:exp_log}

% The experimental track for this thesis evolved substantially as evidence accumulated about what the Five Chambers environment, the Hebbian hypothesis, and the available compute could jointly support. This section records the methodological choices and the rationale behind each pivot, so that the contribution of the final experimental design can be read against the alternatives considered.

% \paragraph{Initial experiment.}
% The first experimental track combined three parts: a MindForge-style LLM agent as the per-agent reasoning core, a multi-agent actor-critic learner (MAPPO with a shared centralised critic, with IPPO as an ablation) that fine-tuned a LoRA adapter on top of the frozen base, and the Hebbian social-plasticity graph. The choice reflected the thesis's position relative to prior work: the Hebbian graph is proposed as an augmentation of existing cooperative MARL methods, not a replacement for them. The experimental grid was sized accordingly: model-axis baselines, communication on/off, Hebbian on/off, and a population-scaling sweep at three, six, and nine agents.


% \paragraph{What the initial experiment revealed.}
% Runs through this stack converged on the same ceiling. Agents plateaued at Chamber 1 across every condition, regardless of whether the reinforcement learner or the Hebbian module was active; the furthest skill-driven milestone was a Chamber 1 pickup task. Subsequent chambers were entered only through a timeout teleport added to prevent indefinite stalls in solo learning. In the runs that combined reinforcement learning with Hebbian shaping, per-agent returns were unequal: a single agent dominated task reward while its teammates trended net-negative. The bond matrix that the Hebbian module was meant to grow instead collapsed, with mean bond strength falling from its warm-start value to near zero. The reason was structural rather than algorithmic: the rule's co-activity input was empty for agents that never co-engaged.

% This made the core hypothesis untestable in this stack. The thesis claim is that reward-modulated Hebbian plasticity over inter-agent bonds builds useful social structure; verifying or refuting that requires agents to co-engage often enough for the rule to fire. With agents stuck in solo skills in Chamber 1, there was no co-activity to reinforce. The aggregated cooperation score, intended as a summary statistic, was correspondingly uninformative, it sat near its noise floor regardless of behaviour and could only meaningfully report deltas, not absolute levels.

% \paragraph{Isolating the Hebbian representation from the gradient.}
% Before moving on, we ran two inference-time-only studies that address a different and methodologically cleaner question: when the language model is not trained at all, does the Hebbian representation alone improve coordination? The first variant surfaces the current bond matrix as a social-context block in the system prompt; the second additionally surfaces the rewards a teammate's actions have just propagated to the agent. Both variants disable the reinforcement-learning gradient entirely. If either condition outperforms the plain language-model baseline, the Hebbian representation has signal that does not depend on a successful gradient path; if both match the baseline, the contribution must flow through the optimiser instead. These runs are inexpensive, isolate a question the gradient-driven experiments cannot, and remain part of the final experimental matrix regardless of how the policy-optimisation track resolves.
